% !PRINTED 672
\weberpage{672}
%<<CONT
das wir durch die Gleichung

\weq{2}{\lambda_s = \delta_s \log|\eta_s|}

definiren, worin $|\eta_s|$ den absoluten Werth von $\eta_s$ und $\log|\eta_s|$ den reellen natürlichen Logarithmus der positiven Zahl $|\eta_s|$ bedeutet. Ist $\Omega_s$ reell und das Zeichen $\pm$ so gewählt, dass $\pm \eta_s$ positiv ist, so ist

\[ \lambda_s = \log(\pm \eta_s). \]

Ist aber $\eta_s, \eta_s'$ ein imaginäres Paar, so ist

\[ \lambda_s = \log \eta_s \eta_s', \]

und zu $\eta_s$ und $\eta_s'$ gehört dasselbe $\lambda_s$, so dass die Anzahl der verschiedenen $\lambda_s$ gleich $\nu$ ist. Aus dieser Bestimmung ergiebt sich noch

\weq{3}{\lambda_1 + \lambda_2 + \cdots + \lambda_\nu = \log N_a(\eta),}

wenn, wie früher, $N_a$ die absolute Norm bedeutet.

Die Zahlen $\lambda_1, \lambda_2, \ldots, \lambda_\nu$ wollen wir die \gesperrt{conjugirten Logarithmen} der Zahl $\eta$ nennen.

Wenn $\eta$ eine \gesperrt{ganze Zahl} des Körpers $\Omega$ ist, so ist die absolute Norm eine \gesperrt{natürliche Zahl}, die nur dann $= 1$ ist, wenn $\eta$ eine Einheit ist, und daraus folgt nach (3):

1. \gesperrt{Die Summe der conjugirten Logarithmen einer ganzen Zahl $\eta$ ist positiv, und nur dann gleich Null, wenn $\eta$ eine Einheit ist.}

Bedeutet $\omega_1, \omega_2, \ldots, \omega_n$ eine Basis der ganzen Zahlen in $\Omega$ (eine Minimalbasis von $\Omega$, \S.\,145), und $\omega_{1,s}, \omega_{2,s}, \ldots, \omega_{n,s}$ für $s = 1, 2, \ldots, n$ die conjugirten Zahlen, so lassen sich alle ganzen Zahlen des Körpers $\Omega_s$ in der Form darstellen:

\weq{4}{\eta_s = x_1 \omega_{1,s} + x_2 \omega_{2,s} + \cdots + x_n \omega_{n,s},}

mit ganzen rationalen Co\"efficienten $x_1, x_2, \ldots, x_n$. Die Determinante dieses Systemes

\[ \Sigma \pm \omega_{1,1} \omega_{2,2} \ldots \omega_{n,n} = \sqrt{\Delta} \]

ist die Quadratwurzel aus der Discriminante $\Delta$ des Körpers, und demnach immer von Null verschieden. Demnach lässt sich das System (4) nach den Unbekannten $x_1, x_2, \ldots, x_n$ auflösen, und wenn wir nun festsetzen, dass die absoluten Werthe der Zahlen $\eta_s$ nicht über eine gewisse Grenze hinausgehen sollen, so ergeben sich aus diesen Auflösungen Grenzen für die ganzen rationalen Zahlen $x_i$.
